% GENERATED by scripts/build-lecture-slides.py; do not edit. Portable profile: XeTeX, theme and fonts from the core.
\makeatletter\IfFileExists{core/beamer/beamerthememx.sty}{\def\mxcorepath{core/}\def\input@path{{core/beamer/}{core/}}}{\def\mxcorepath{../core/}\def\input@path{{../core/beamer/}{../core/}}}\makeatother
\documentclass[aspectratio=169,10pt,t]{beamer}
\usetheme{mx}
\usepackage{graphicx,booktabs}
\usepackage[french]{babel}
\setlength{\parskip}{.45em}
\hypersetup{pdfauthor={},pdfsubject={Lecture slides},pdfkeywords={},pdftitle={Branches infinies et asymptotes}}
\providecommand{\N}{\mathbb{N}}
\providecommand{\Z}{\mathbb{Z}}
\providecommand{\D}{\mathbb{D}}
\providecommand{\Q}{\mathbb{Q}}
\providecommand{\R}{\mathbb{R}}
\providecommand{\C}{\mathbb{C}}
\providecommand{\vect}[1]{\overrightarrow{#1}}
\providecommand{\vi}{\vec{\imath}}
\providecommand{\vj}{\vec{\jmath}}
\providecommand{\vk}{\vec{k}}
\providecommand{\norm}[1]{\left\lVert #1\right\rVert}
\providecommand{\vc}[2]{\begin{pmatrix}#1\\#2\end{pmatrix}}
\providecommand{\deter}[2]{\det\left(#1,#2\right)}
\providecommand{\abs}[1]{\left\lvert #1\right\rvert}
\providecommand{\dd}{\,\mathrm{d}}
\providecommand{\e}{\mathrm{e}}
\providecommand{\Card}{\operatorname{Card}}
\providecommand{\E}{\operatorname{E}}
\providecommand{\V}{\operatorname{V}}
\providecommand{\degre}{^{\circ}}
\providecommand{\un}[1]{\,\mathrm{#1}}
\providecommand{\sci}[2]{#1\times10^{#2}}
\providecommand{\pgcd}{\operatorname{PGCD}}
\providecommand{\ppcm}{\operatorname{PPCM}}
\title{Branches infinies et asymptotes}
\subtitle{Cours et exemples}
\date{}
\author{}
\begin{document}
\mxtitleframe
\begin{frame}{Objectifs}
Reconnaître une asymptote verticale, horizontale ou oblique et étudier la position relative d’une courbe et de sa droite asymptote.
\end{frame}
\begin{frame}{Une courbe qui se rapproche d’une droite}
La droite \(x=a\) est une asymptote verticale si au moins une limite latérale de \(f\) en \(a\) vaut \(+\infty\) ou \(-\infty\). La droite \(y=\ell\) est une asymptote horizontale en \(+\infty\) (respectivement \(-\infty\)) si \(\lim_{x\to+\infty}f(x)=\ell\) (respectivement en \(-\infty\)).
\end{frame}
\begin{frame}{Une courbe qui se rapproche d’une droite}
Pour \(f(x)=2+\dfrac1{x-1}\), la droite \(x=1\) est asymptote verticale et \(y=2\) est asymptote horizontale aux deux infinis. \(f(x)-2=1/(x-1)\) donne une courbe sous \(y=2\) à gauche de \(1\), au-dessus à droite.
\end{frame}
\begin{frame}{Asymptote oblique}
La droite \(y=ax+b\), avec \(a\ne0\), est une asymptote oblique en l’infini considéré si \(f(x)-(ax+b)\) tend vers \(0\).
\end{frame}
\begin{frame}{Asymptote oblique}
Si \(\lim f(x)/x=a\in\mathbb{R}\) et \(\lim(f(x)-ax)=b\in\mathbb{R}\), alors \(y=ax+b\) est une asymptote; elle est oblique si \(a\ne0\), horizontale si \(a=0\).
\end{frame}
\begin{frame}{Asymptote oblique}
La seconde limite signifie exactement \(f(x)-(ax+b)\to0\). Réciproquement, poser \(r(x)=f(x)-(ax+b)\) donne \(r(x)\to0\) et \(f(x)/x=a+b/x+r(x)/x\to a\).
\end{frame}
\begin{frame}{Asymptote oblique}
Pour une fraction rationnelle dont le numérateur a un degré de plus que le dénominateur, effectuer la division. Si \(f(x)=ax+b+r(x)\) avec \(r(x)\to0\), la droite cherchée est \(y=ax+b\). Étudier le signe de \(r\) pour situer la courbe.
\end{frame}
\begin{frame}{Asymptote oblique}
\[
\dfrac{x^2+1}{x-1}=x+1+\dfrac2{x-1}\qquad(x\ne1).
\]
L’asymptote oblique aux deux infinis est \(y=x+1\). La courbe est en dessous pour \(x<1\), au-dessus pour \(x>1\). \(x=1\) est également une asymptote verticale.
\end{frame}
\begin{frame}{Toutes les branches n’ont pas d’asymptote}
Si \(f(x)\to\pm\infty\) et \(|f(x)/x|\to+\infty\), on parle d’une branche parabolique de direction l’axe des ordonnées. Si \(f(x)\to\pm\infty\) et \(f(x)/x\to0\), la direction est l’axe des abscisses. Si \(f(x)/x\to a\ne0\) mais \(f(x)-ax\to\pm\infty\), la branche a pour direction \(y=ax\), sans asymptote oblique. Ce vocabulaire décrit une direction, pas nécessairement une parabole.
\end{frame}
\begin{frame}{Toutes les branches n’ont pas d’asymptote}
Soit \(F(x)=(x^2+3x+5)/(x+1)\), pour \(x\ne-1\). Comme
\[
(x+1)(x+2)=x^2+3x+2,
\qquad F(x)=x+2+\dfrac3{x+1},
\]
la droite \(y=x+2\) est asymptote aux deux infinis. L’écart vaut \(3/(x+1)\) : la courbe est sous la droite sur \(]-\infty;-1[\) et au-dessus sur \(]-1;+\infty[\). Au voisinage de \(-1\), le numérateur du reste reste positif; les limites latérales sont donc \(-\infty\) et \(+\infty\). La verticale est \(x=-1\).
\end{frame}
\begin{frame}{Toutes les branches n’ont pas d’asymptote}
Dans chaque direction infinie, commencer par \(\lim f\). Si elle est finie, on a une horizontale. Si elle est infinie, calculer \(\lim f/x\). Une valeur finie \(a\ne0\) donne une direction possible, mais il reste à calculer \(\lim(f-ax)\) : une valeur finie donne l’ordonnée à l’origine de l’asymptote; une valeur infinie ne donne pas de droite asymptote. Traiter \(+\infty\) et \(-\infty\) séparément.
\end{frame}
\begin{frame}{Toutes les branches n’ont pas d’asymptote}
Pour \(H(x)=\sqrt{x^2+4}\), la racine dominante est \(|x|\), pas \(x\) partout. En \(+\infty\),
\[
H(x)-x=\dfrac4{\sqrt{x^2+4}+x}\longrightarrow0.
\]
En \(-\infty\), c’est \(H(x)+x=4/(\sqrt{x^2+4}-x)\) qui tend vers \(0\). Les deux asymptotes sont donc \(y=x\) à droite et \(y=-x\) à gauche. Les écarts sont positifs; la courbe reste au-dessus de chacune dans la comparaison correspondante.
\end{frame}
\begin{frame}{Toutes les branches n’ont pas d’asymptote}
La définition concerne le comportement à l’infini. Elle n’interdit pas des intersections à des abscisses finies. Pour en trouver, résoudre l’équation de l’écart égal à zéro sur le domaine, au lieu d’imposer que la courbe reste toujours du même côté.
\end{frame}
\end{document}
